\begin{explanationbox}
	\textbf{\textcolor{CExplanation}{一句话直觉}}

	抛物线$y^2=2px$上点$(x_0,y_0)$处的切线方程$y_0y=p(x+x_0)$可以看作将原方程中的$y^2$替换为$y_0y$，将$2px$替换为$p(x+x_0)$，本质是将二次项线性化。

	\medskip
	\textbf{\textcolor{CExplanation}{核心拆解}}
	\begin{description}[style=nextline, leftmargin=2.8em, labelsep=0.5em, itemsep=0.45em]
		\item[\textbf{要点一（斜率推导）：}] 对$y^2=2px$两边求导得$2y y'=2p$，故斜率$k=\frac{p}{y_0}$（$y_0\neq0$）。当$y_0=0$时，切线为$x=0$。
		\item[\textbf{要点二（方程化简）：}] 切线方程$y-y_0=k(x-x_0)$代入$k$并利用$y_0^2=2px_0$化简为$y_0y=p(x+x_0)$。
	\end{description}

	\medskip
	\textbf{\textcolor{CExplanation}{几何本质（点线唯一交点）}}

	切线与抛物线只有一个公共点$P$，且$P$在抛物线上，这正是切线的定义。

	\medskip
	\textbf{\textcolor{CExplanation}{代数意义（代换法则）}}

	切线方程可由原方程通过“半代换”得到：将$y^2$换成$y_0y$，将$x$换成$\frac{x+x_0}{2}$，系数$2p$对应$p$。这种代换对一般二次曲线也适用。

	\medskip
	\textbf{\textcolor{CExplanation}{考点价值（常见考法）}}
	\begin{itemize}[leftmargin=*, itemsep=0.5em, topsep=0.4em]
		\item \textbf{考法一（直接写切线）：}已知切点坐标$(x_0,y_0)$和抛物线方程，直接代入公式得切线方程。
		\item \textbf{考法二（求切点坐标）：}已知切线经过某点或斜率，设切点$(x_0,y_0)$，利用切线公式和点在抛物线上联立求解。
	\end{itemize}

	\medskip
	\textbf{\textcolor{CExplanation}{顿悟点}}

	\textbf{公式的对称性：}$y_0y$和$x+x_0$体现了点$(x_0,y_0)$的“一半”替换。当$y_0=0$时公式仍给出$x=0$，说明顶点处的切线为垂直于$x$轴的直线。

	\medskip
	\textbf{\textcolor{CExplanation}{使用场景}}
	\begin{itemize}[leftmargin=*, itemsep=0.5em, topsep=0.4em]
		\item \textbf{场景一（切点已知）：}直接代公式，快速写出切线，不需要判别式法。
		\item \textbf{场景二（切线已知）：}已知某直线与抛物线相切，可设切点并用切线方程表示直线，联立条件求解。
	\end{itemize}
\end{explanationbox}
